\documentclass[a4paper]{article}     
\usepackage{amsfonts}
\usepackage{amsmath}
\usepackage{amssymb}
\usepackage{graphicx}

\newcommand{\integral}{\displaystyle\int}

\begin{document}

\noindent \large Auteur: Elias Alsa Ati         \\
\noindent \large Studierichting: Informatica    \\
\noindent \large Oefeningenreeks 4A nr. 43.1    \\

\medskip

\normalsize

$$\integral \dfrac{dx}{x(x-1)^2}$$

We splitsen de integraal op in partiële Breuken: \\

$\begin{array}{l}
    
    \dfrac{1}{x(x-1)^2} \Rightarrow \dfrac{A}{x} + \dfrac{B}{(x-1)^2} + \dfrac{C}{(x-1)} \\
    
    \left.
    \begin{array}{l}
        A = \dfrac{1}{(x-1)^2}  \\
    \end{array}
    \right|_{x=0} = 1           \\

    \left.
    \begin{array}{l}
        B = \dfrac{1}{x}        \\
    \end{array}
    \right|_{x=1} = 1           \\

    \left.
    \begin{array}{ll}
        C   &= \left(\dfrac{1}{x(x-2)^2} - \dfrac{B}{(x-1)^2}\right) \cdot (x - 1)     \\
            &= \left(\dfrac{1 - x}{x(x-1)^2}\right) \cdot (x-1)         \\
            &= \left(\dfrac{-(x-1)^2}{x(x-1)^2}\right) = - \dfrac{1}{x} \\
    \end{array}
    \right|_{x=1} = -1                                                  \\
\end{array}$    \\

Hier uit weten we nu dat $A = 1$, $B = 1$, $C = -1$ \\

en we vullen het in de opsplitsing:                 \\

$$ \Rightarrow \integral \dfrac{A}{x} dx + \integral \dfrac{B}{(x-1)^2} dx + \integral \dfrac{C}{x-1} dx $$

$\begin{array}{l}
    = \integral \dfrac{1}{x} dx + \integral \dfrac{1}{(x-1)^2} dx - \integral \dfrac{1}{x-1} dx \\
    = \ln|x| + \integral \dfrac{1}{(x-1)^2} dy - \integral \dfrac{1}{x-1} dx    \\
    = \ln|x| + \integral y^{-2} dy - \integral \dfrac{1}{y} dy    \\
    = \ln|x| - \dfrac{1}{y} - \ln|y| + c    \\
    = \ln|x| - \dfrac{1}{x-1} - \ln|x-1| + c    \\
\end{array}
\begin{array}{l}
    y = x - 1 \Rightarrow x = y + 1 \\
    dy = dx     \\
\end{array}$    \\

Je resultaat is dan:    \\

$$\ln|x| - \dfrac{1}{x-1} - \ln|x-1| + c$$

\end{document} 